%======================================================================
\section{Forced branching: the lamp sources}\label{sec:rates}
%======================================================================

This section and the next build source groups with prescribed polynomial profiles. The mechanism is branching forced
at each magnification. A relation $xyx^{-1}=y^m$ makes the spectral law of $y$ nearly invariant under multiplication of the
angle by $m$, and a finite set of involutions permuted cyclically by $y$ makes it nearly invariant under a rotation by
$1/q$ of a turn. A law invariant under both maps puts mass at most $q^{-n}$ on an arc of length $q^{-1}m^{-n+1}$, so it satisfies
a Frostman bound of exponent $s=\log_mq$ and has dimension at least $s$; Haar measure, for instance, is invariant and has dimension
one. A spectral law that is nearly invariant is close to such a law, so it needs about $e^{-s}$ atoms and $s\log(1/e)$ bits of
entropy.

\paragraph{Turns.} In this section and the next we identify $\Tor$ with $\R/\Z$ through $\vartheta=2\pi t$ and measure
lengths and Wasserstein distances in turns: $d(t,u)\in[0,\frac12]$, and $\Wone$ is $1/(2\pi)$ times the distance of
\cref{sec:transport}. \Cref{lem:spectral-coupling} then reads $\Wone(\mu^\rho_P,\mu^\rho_Q)\le\frac14\hs{P\xi-\xi Q}$ and
$\Wone(\mu^{U^*\rho U}_Y,\mu^\rho_Y)\le\frac12\hs{[U,\xi]}$, and a $(\kappa,s)$-spread law satisfies $\nu(I)\le(\kappa|I|)^s$. Put
$D(t)=mt$ and $S(t)=t+1/q$ modulo $1$.

\subsection{The sources}\label{sec:rates-sources}

\begin{definition}[Lamp sources]\label{def:lamps-source}
For integers $q\ge1$ and $m\ge2$, the \emph{lamp source} $L_{q,m}$ has generators $z_0,\dots,z_{q-1},w,y,x$ and the relators
\begin{gather*}
z_j^2,\quad w^2,\quad[z_j,z_k],\quad[z_j,w],\quad z_0z_1\cdots z_{q-1}w,\\
yz_jy^{-1}z_{j+1}^{-1}\ (j\bmod q),\qquad xwx^{-1}w^{-1},\qquad xyx^{-1}y^{-m},
\end{gather*}
with distinguished involution $w$. The row makes $w$ commute with the lamps $z_j$, and conjugating it by $y$ permutes its
factors cyclically, so $w$ is central. Its profiles are those of \cref{ex:presented}(a): a model at precision $e$ is a
family of unitaries $Z_j,W,Y,X$ and a density matrix $\rho$, $\xi=\rho^{1/2}$, with every relator of defect at most $e$ in
the state, every generator of state commutator at most $e$, and $\hs{(W+I)\xi}\le e$.
\end{definition}

For $q=2$ the row gives $z_1=z_0w$, and $L_{2,2}$ is the group $K$ of Slofstra and Vidick~\eqref{eq:K}, with $a,b,c$ the
two lamps and $w$. The power relation is a dilation of the kind of Slofstra's HNN constructions~\cite[Section~5]{Slo19}
and of Fu's self-tests~\cite{Fu22}.

\begin{theorem}[One lamp source]\label{thm:lamp}
Let $q=2^r$ with $r\ge1$, let $q$ divide $m$, and put $s=\log_mq$. There are constants $c,C>0$ such that every model of
$L_{q,m}$ at precision $e\le\frac12$ on $\C^d$, with density matrix $\rho$, has
\[
d\ \ge\ \rank\rho\ \ge\ c\,e^{-s}\qquad\text{and}\qquad S(\rho)\ \ge\ s\log(1/e)-C ,
\]
and for every $e\le\frac12$ there is a model at precision $e$ with the flat state, of dimension at most $Ce^{-s}$ and
entropy at most $s\log(1/e)+C$. So $h_{L_{q,m}}(e)=\Theta(e^{-s})$ and $h^S_{L_{q,m}}(e)=s\log(1/e)+O(1)$.
\end{theorem}

The proof occupies the rest of the section. Its steps are a rotation implemented on the whole state
(\cref{sec:rates-rotation}), an invariant law nearby (\cref{sec:rates-shadow}), the spread of every invariant law
(\cref{sec:rates-spread}), and matching models (\cref{sec:rates-models}).

\subsection{A rotation on the whole state}\label{sec:rates-rotation}

Work first in the group algebra of $L_{q,m}$, before $w=-1$ is imposed. For a sign vector $\chi\in\{\pm1\}^q$ put
$P_\chi=2^{-q}\prod_j(1+\chi_jz_j)$. These are commuting self-adjoint idempotents with $P_\chi P_{\chi'}=0$ for
$\chi\ne\chi'$, $\sum_\chi P_\chi=1$ and $wP_\chi=(\prod_j\chi_j)P_\chi$, so
\[
\sum_{\chi\in\Omega_-}P_\chi=\tfrac12(1-w),\qquad\Omega_-=\Bigl\{\chi:\prod_j\chi_j=-1\Bigr\}.
\]
With $(\sigma\chi)_k=\chi_{k-1}$, the conjugacies give $yP_\chi y^{-1}=P_{\sigma\chi}$. Every $\chi\in\Omega_-$ has an orbit of
length exactly $q$: if its least period $l$ is smaller, its first $l$ entries repeat $q/l$ times, an even number when $q$ is
a power of two, and the product of all entries is $+1$. Choose a representative $\chi_O$ of each orbit $O\subseteq\Omega_-$,
put $\zeta=e^{2\pi i/q}$, and
\[
T=\sum_O\sum_{a=0}^{q-1}\zeta^aP_{\sigma^a\chi_O}.
\]

\begin{lemma}\label{lem:rotation-selector}
In the group algebra, $T^*T=TT^*=\frac12(1-w)$ and $Ty=\zeta yT$. Consequently, for a model of $L_{q,m}$ at precision $e$
there is a unitary $U$ with
\[
\hs{[U,\xi]}\le Ce\qquad\text{and}\qquad\hs{(UYU^*-\zeta Y)\xi}\le Ce ,
\]
with $C$ depending only on $q$ and $m$.
\end{lemma}

\begin{proof}
By orthogonality, $T^*T=TT^*=\sum_O\sum_aP_{\sigma^a\chi_O}=\frac12(1-w)$. Next $yTy^{-1}=\sum_O\sum_a\zeta^aP_{\sigma^{a+1}\chi_O}=
\zeta^{-1}T$, since $\sigma^q\chi_O=\chi_O$; so $Ty=\zeta yT$. These identities have finite derivations from the defining
relators (expand the products, reduce squares, commute lamps, use the row for $w$, and the $q$ conjugacies), so in the
algebra $L_{q,m}^-$ of \cref{ex:presented}(a) the elements $T^*T-1$, $TT^*-1$ and $Ty-\zeta yT$ lie in the defining ideal.
Let $B=T(\varphi)$. By \cref{lem:finite-identity}, $\hs{(B^*B-I)\xi}$, $\hs{(BB^*-I)\xi}$, $\hs{(BY-\zeta YB)\xi}$ and
$\hs{[B,\xi]}$ are $O(e)$, and $\opn B=O(1)$. \Cref{lem:polar} gives a unitary $U$ with $\hs{(U-B)\xi}$, $\hs{\xi(U-B)}$ and
$\hs{[U,\xi]}$ of order $e$. Then $(UY-\zeta YU)\xi=(BY-\zeta YB)\xi+(U-B)\xi Y+(U-B)[Y,\xi]-\zeta Y(U-B)\xi$ is $O(e)$, and
$(UYU^*-\zeta Y)\xi=(UY-\zeta YU)U^*\xi$ is $O(e)$ after moving $U^*$ past $\xi$.
\end{proof}

No sector is selected: the polynomial $T$ combines all sign vectors of odd parity, and the state is never changed. For
$q=2$, $T=P_{(-,+)}-P_{(+,-)}$, which in the negative quotient is the lamp $z_0$ up to sign, and $U$ is close to it.

\begin{lemma}[Two approximate invariances]\label{lem:invariances}
For a model of $L_{q,m}$ at precision $e$, the spectral law $\mu=\mu^\rho_Y$ satisfies
\[
\delta_D=\Wone(D_*\mu,\mu)\le\tfrac{m+3}2\,e,\qquad\delta_S=\Wone(S_*\mu,\mu)\le C'e .
\]
\end{lemma}

\begin{proof}
Put $Q=XYX^*$ and $R=Y^m$. The relator gives $\hs{(QR^*-I)\xi}\le e$, so $\hs{(Q-R)\xi}\le e+2\hs{[R,\xi]}\le(2m+1)e$; with
$\hs{[Q,\xi]}\le3e$, $\hs{R\xi-\xi Q}\le(2m+4)e$. By~\eqref{eq:spectral-push} and \cref{lem:spectral-coupling} in turns,
$\Wone(D_*\mu,\mu^{X^*\rho X}_Y)\le\frac{m+2}2e$ and $\Wone(\mu^{X^*\rho X}_Y,\mu)\le\frac e2$. For the rotation let $U$ be as in
\cref{lem:rotation-selector}: $\mu^\rho_{\zeta Y}=S_*\mu$, $\hs{\zeta Y\xi-\xi UYU^*}\le\hs{(\zeta Y-UYU^*)\xi}+\hs{[UYU^*,\xi]}=
O(e)$, and $\mu^\rho_{UYU^*}=\mu^{U^*\rho U}_Y$ is within $\frac12\hs{[U,\xi]}$ of $\mu$.
\end{proof}

\subsection{An invariant law nearby}\label{sec:rates-shadow}

For $m=2$ and $q=2$ the invariant law is unique, and \cref{lem:circle} gives a direct estimate. In general there are many
invariant laws, Cantor-like ones included, and we show instead that a nearly invariant law is close to \emph{some} invariant
law, by shadowing a stationary pseudo-orbit.

\begin{lemma}[Shadowing]\label{lem:shadowing}
Let $q$ divide $m$, and let $\mu$ be a probability measure on $\R/\Z$ with $\delta_D=\Wone(D_*\mu,\mu)$ and
$\delta_S=\Wone(S_*\mu,\mu)$. There is a probability measure $\nu$ with $D_*\nu=\nu$ and $S_*\nu=\nu$ such that
\[
\Wone(\mu,\nu)\le b+\frac{\delta_D+b}{m-1},\qquad b=\frac{(q-1)\delta_S}2 .
\]
\end{lemma}

\begin{proof}
\emph{Symmetrize.} Let $\mu_0=\frac1q\sum_{k<q}S^k_*\mu$. Since $S$ is an isometry, $\Wone(\mu,S^k_*\mu)\le k\delta_S$, so
$\Wone(\mu,\mu_0)\le b$. As $q$ divides $m$, $DS=D$, so $D_*\mu_0=D_*\mu$ and $\delta:=\Wone(D_*\mu_0,\mu_0)\le\delta_D+b$.

\emph{A symmetric coupling.} Let $\eta$ be an optimal coupling of $D_*\mu_0$ and $\mu_0$. Disintegrate $\mu_0$ along the fibres
of $D$; since $S$ preserves every fibre and $\mu_0$, the conditional laws can be replaced by their averages over the powers of
$S$. Choosing $(z,v)$ from $\eta$ and then $u$ from the conditional law of $\mu_0$ on the fibre $D^{-1}(z)$ gives a law $Q$ of
$(u,v)$ with both marginals $\mu_0$, with $Du=z$, so that $\int d(Du,v)\,dQ=\delta$, and invariant under $(u,v)\mapsto(Su,v)$.

\emph{A stationary chain.} Let $(U_t)_{t\in\Z}$ be the stationary Markov chain with stationary law $\mu_0$ whose consecutive
pairs $(U_t,U_{t+1})$ have law $Q$. The conditional law of $U_t$ given the whole future $(U_{t+1},U_{t+2},\dots)$ is the
reverse kernel of $Q$ at $U_{t+1}$, by the Markov property of the reversed chain; by the symmetry of $Q$ it is invariant
under $S$.

\emph{Shadowing.} Let $E_t\in[-\frac12,\frac12)$ represent $mU_t-U_{t+1}$ modulo $1$, so $\mathbb E|E_t|=\delta$, and put
\[
V_t=U_t-\sum_{j\ge0}m^{-j-1}E_{t+j}\pmod1 .
\]
The series converges absolutely, and $mV_t-V_{t+1}=mU_t-E_t-U_{t+1}=0$ modulo $1$. The process $(V_t)$ is a fixed function of
the shifted process, so it is stationary, and the law $\nu$ of $V_0$ satisfies $D_*\nu=\nu$. Also
$\mathbb E\,d(U_0,V_0)\le\sum_jm^{-j-1}\mathbb E|E_j|=\delta/(m-1)$, so $\Wone(\mu_0,\nu)\le\delta/(m-1)$.

\emph{Symmetry.} $E_0$ depends on $U_0$ only through $mU_0$, which is unchanged when $U_0$ is replaced by $SU_0$ because $q$
divides $m$, and every later $E_j$ depends only on the future. So $V_0=U_0-g(mU_0,U_1,U_2,\dots)$ for a measurable $g$, and
replacing $U_0$ by $SU_0$ replaces $V_0$ by $SV_0$. Conditional on the future, $U_0$ has an $S$-invariant law; hence so has
$V_0$, and $\nu$ is $S$-invariant. Finally $\Wone(\mu,\nu)\le\Wone(\mu,\mu_0)+\Wone(\mu_0,\nu)$.
\end{proof}

No uniqueness, ergodicity or continuity of an inverse branch is used, and the pseudo-orbit errors are controlled only on
average.

\subsection{Every invariant law is spread}\label{sec:rates-spread}

\begin{lemma}[Spread]\label{lem:spread}
Let $q\ge2$ divide $m$ and $s=\log_mq$. Every probability measure $\nu$ on $\R/\Z$ with $D_*\nu=\nu$ and $S_*\nu=\nu$ has no
atoms and satisfies $\nu(I)\le(q|I|)^s$ for every arc $I$; it is $(q,s)$-spread. Consequently every probability measure
$\lambda$ with at most $d$ atoms has $\Wone(\lambda,\nu)\ge d^{-1/s}/(4qm)$.
\end{lemma}

\begin{proof}
For a point $t$, the $q$ points $S^kt$ are distinct and all map to $Dt$, so $q\nu(\{t\})\le\nu(\{Dt\})$, and iterating gives
$\nu(\{t\})\le q^{-n}$ for every $n$. For a half-open arc $I$ of length at most $1/q$, the arcs $S^kI$ are disjoint and lie in
$D^{-1}(DI)$, so $q\nu(I)\le\nu(D^{-1}(DI))=\nu(DI)$, where $DI$ is an arc of length $\min(m|I|,1)$. For $0<|I|\le1/q$ let
$n=\lfloor\log_m(1/(q|I|))\rfloor+1$; the first $n$ images of $I$ have length at most $1/q$ before each step, so
$\nu(I)\le q^{-n}<(q|I|)^s$. Longer arcs satisfy the bound trivially, and endpoints do not matter.

Let $F$ be the set of atoms of $\lambda$, $r=(2d)^{-1/s}/(2q)$, and $f(t)=d(t,F)$. The points within $r$ of $F$ form at most
$d$ arcs of length $2r$, of total $\nu$-mass at most $d(2qr)^s=\frac12$. So $\int f\,d\nu\ge r/2$ while $\int f\,d\lambda=0$,
and $\Wone(\lambda,\nu)\ge r/2=2^{-1/s}d^{-1/s}/(4q)\ge d^{-1/s}/(4qm)$, since $2^{1/s}=m^{\log_q2}\le m$.
\end{proof}

\begin{proof}[Proof of the lower bounds in \cref{thm:lamp}]
By \cref{lem:invariances,lem:shadowing} there is an invariant $\nu$ with $\Wone(\mu,\nu)\le C''e$. The law $\mu$ has at most
$d$ atoms, so \cref{lem:spread} gives $d\ge(4qmC''e)^{-s}$. In radians $\nu$ is $(q,s)$-spread and $\Wone(\mu,\nu)\le2\pi C''e$,
and $\hs{[Y,\xi]}\le e$, so \cref{thm:entropy-transport} gives $S(\rho)\ge s\log(1/e)-C$. Finally $\rank\rho\ge2^{S(\rho)}$.
\end{proof}

The entropy bound uses the state commutator of $Y$, not only the spread of its spectral law: by \cref{ex:clock} a pure
state can have a spread spectral law.

\subsection{Matching models}\label{sec:rates-models}

\begin{example}[Fourier lamps]\label{ex:fourier-lamps}
Let $q$ divide $m$, $L\ge1$, $h=m^L-1$, and index a basis of $\C^{q^L}$ by words $b=(b_0,\dots,b_{L-1})$ over
$\{0,\dots,q-1\}$. Put
\[
t_b=\frac{m}{qh}\sum_{j=0}^{L-1}b_jm^{L-1-j},\qquad Y=\mathrm{diag}(e^{2\pi it_b}),\qquad W=-I,
\]
let $X$ be the permutation matrix with $X_{b,Rb}=1$, where $R$ rotates a word one place to the left, and on the first letter
let $v_j(a)=q^{-1/2}\zeta^{ja}$, $E_j=|v_j\rangle\langle v_j|$ and $Z_j=(I_q-2E_j)\otimes I_{q^{L-1}}$. Then all relators of
$L_{q,m}$ hold exactly except the lamp conjugacies, whose defects are at most $4\pi(q-1)/(q(m^L-1))$ in operator norm.
\end{example}

\begin{proof}
The $E_j$ are orthogonal rank-one projections with sum $I_q$, so the $Z_j$ are commuting involutions with
$\prod_jZ_j=I-2\sum_jE_j=-I$, and the row, the squares and the commutators hold; $W=-I$ is central. An integer calculation
gives $mt_b-t_{Rb}=(m/q)b_0\in\Z$, so $XYX^*=Y^m$ exactly. Write $t_b=b_0/q+b_0/(qh)+r_b$, with $r_b$ depending only on
$b_1,\dots,b_{L-1}$, and let $Y_0=D_0\otimes\mathrm{diag}(e^{2\pi ir_b})$ with $D_0=\mathrm{diag}(\zeta^{a})$ on the first letter. Then
$D_0E_jD_0^*=E_{j+1}$, so $Y_0Z_jY_0^*=Z_{j+1}$, including $j=q-1$, and $\opn{Y-Y_0}\le\max_{b_0}|e^{2\pi ib_0/(qh)}-1|\le
2\pi(q-1)/(qh)$. Hence $\opn{YZ_jY^*-Z_{j+1}}\le2\opn{Y-Y_0}$, and this is the defect of the relator, since its last factor is
unitary.
\end{proof}

\begin{proof}[Proof of the upper bounds in \cref{thm:lamp}]
With the flat state every state commutator vanishes and every defect in the state is at most its operator norm. Choose
$L$ least with $4\pi(q-1)/(q(m^L-1))\le e$; then $m^{L-1}\le C/e$, and the dimension $q^L=q\,(m^{L-1})^s\le Ce^{-s}$ has
entropy $L\log q\le s\log(1/e)+C$.
\end{proof}

\begin{remark}[The arithmetic of the lamp count]\label{rem:lamp-arithmetic}
Write $q=2^vt$ with $t$ odd. Every odd-parity sign vector has an orbit whose length $l$ satisfies that $q/l$ is odd, so $2^v$
divides $l$. Grading each orbit by powers of $e^{2\pi i/2^v}$ gives, by the same proof, a rotation by $2^{-v}$ of a turn on the
whole state; if $2^v$ divides $m$, the exponent of $L_{q,m}$ is $\log_m2^v$, with matching models obtained from $2^v$ Fourier
lamps, each repeated $t$ times. For odd $q$ there is an exact one-dimensional model, with every lamp and $w$ equal to $-1$ and
$x=y=1$, and the profile is bounded. So the power of two in $q$ is what branches; an odd factor adds nothing. The
divisibility of $m$ by the rotation order is used in $DS=D$, in the shadowing and in the models.
\end{remark}

\subsection{A five-generator source with the same profile}\label{sec:rates-five}

The lamps can be replaced by a single central element of order $q$. For $q=2^r$ dividing $m$, let
$G_{q,m}=\langle a,c,w,x,y\rangle$ with the relators $a^q$, $c^q$, $w^2$, the commutators of $c$ with $a,x,y,w$ and of $w$ with
$a,x,y$, the relation $w=c^{q/2}$, and $aya^{-1}=cy$, $xyx^{-1}=y^m$, the marked involution being $w$. In the group algebra
put $P_j=q^{-1}\sum_l\zeta^{-jl}c^l$; these are orthogonal projections commuting with $a,x,y,w$, with $cP_j=\zeta^jP_j$ and
$\sum_{j\text{ odd}}P_j=\frac12(1-c^{q/2})$. For odd $j$ let $k_j$ be the inverse of $j$ modulo $q$ and put
$T=\sum_{j\text{ odd}}P_ja^{k_j}$. Then $T^*T=TT^*=\frac12(1-w)$ and $Ty=\zeta yT$ in the full group algebra, by finite derivations
from the relators, so the conclusion of \cref{lem:rotation-selector} holds for $G_{q,m}$, and with it every argument of this
section. Hence $G_{q,m}$ has minimum ambient dimension and minimum rank $\Theta(e^{-s})$ and minimum entropy
$s\log(1/e)+O(1)$, $s=\log_mq$; the abandoned variant with $q=6$ and $c=-1$ does not create a six-way exponent, since its
negative sector can be two-dimensional. We use the lamp sources below because they lie in the restricted linear input
class.

\subsection{Additivity over copies, without a joint spectral measure}\label{sec:additive}

For $i=1,\dots,k$ fix $q_i=2^{r_i}$ dividing $m_i$ and a copy of $L_{q_i,m_i}$, identify the marked involutions, and add to the
relators every commutator between generators of different copies. Call this finite presentation $L_{\mathbf q,\mathbf m}$, and put
\[
s_i=\log_{m_i}q_i,\qquad s=\sum_is_i .
\]

\begin{theorem}[Additive profiles]\label{thm:additive}
There are constants $c,C>0$, depending on the fixed presentation, such that every model of $L_{\mathbf q,\mathbf m}$ at precision
$e\le\frac12$, with density matrix $\rho$ on $\C^d$, satisfies
\[
S(\rho)\ \ge\ s\log(1/e)-C,\qquad d\ \ge\ \rank\rho\ \ge\ c\,e^{-s},
\]
and tensor products of the Fourier-lamp models of \cref{ex:fourier-lamps} give models of dimension at most $Ce^{-s}$ and
entropy at most $s\log(1/e)+C$. So the minimum dimension and the minimum rank are $\Theta(e^{-s})$, and the minimum entropy is
$s\log(1/e)+O(1)$.
\end{theorem}

The difficulty is that the matrices $Y_i$ of different copies commute only approximately, in a weighted norm, and
$\rho$ may be singular or far from flat. We do not assign them a joint spectral measure, and we do not round them to
commuting matrices. Instead we measure spectral arcs of short words copy by copy, in reverse order of depth, and show
that the record has additive entropy and that Alice can predict it at bounded cost.

\emph{Notation.} Norms are unnormalized Hilbert--Schmidt norms, $\hs\xi=1$, and $c(V)=\hs{[V,\xi]}$. Constants $C$ depend on the
fixed sources and on $k$, never on $e$, the depth, the dimension or $\rho$. Spectral angles are in turns.

\begin{lemma}[Short words]\label{lem:short-words}
For each copy there is a unitary $U_i$ with $c(U_i)\le Ce$, $\hs{(U_iY_iU_i^*-\zeta_iY_i)\xi}\le Ce$, $\zeta_i=e^{2\pi i/q_i}$, and
$\hs{[U_i,V]\xi}\le Ce$ for every generator $V$ of another copy. Put $Z_{i,j}=X_i^jY_iX_i^{-j}$ and $T_{i,j}=X_i^jU_iX_i^{-j}$. For
depths $j,l\le N$,
\begin{gather*}
c(Z_{i,j})+c(T_{i,j})\le C(N+1)e,\\
\hs{[Z_{i,j},Z_{h,l}]\xi}+\hs{[T_{i,j},Z_{h,l}]\xi}\le C(N+1)^3e\quad(i\ne h),\\
\hs{[Z_{i,j},Z_{i,l}]\xi}\le Ce\bigl(m_i^{|l-j|}+N+1\bigr),\qquad\hs{(T_{i,j}Z_{i,j}T_{i,j}^*-\zeta_iZ_{i,j})\xi}\le C(N+1)e,\\
\hs{[T_{i,j},Z_{i,l}]\xi}\le Ce\bigl(m_i^{l-j}+N+1\bigr)\quad(l>j).
\end{gather*}
\end{lemma}

\begin{proof}
The unitary $U_i$ is that of \cref{lem:rotation-selector}. The selector polynomial $B_i$ uses only copy-$i$ generators, so its
commutator with a generator $V$ of another copy has a finite certificate and is $O(e)$ on $\xi$; and $[U_i-B_i,V]\xi$ is
$O(e)$ because both one-sided differences $(U_i-B_i)\xi$ and $\xi(U_i-B_i)$ are $O(e)$ and $c(V)\le e$.

For unitaries $R,S$, left multiplication by $R$ and right multiplication by $S$ are commuting unitaries on the
Hilbert--Schmidt space, so
\begin{equation}\label{eq:left-right-power}
\hs{R^a\xi-\xi S^a}\le a\hs{R\xi-\xi S}\qquad(a\ge0).
\end{equation}
Let $b_d=\hs{(X^dYX^{-d}-Y^{m^d})\xi}$ for one copy. Then $X^{d+1}YX^{-d-1}-Y^{m^{d+1}}=X(X^dYX^{-d}-Y^{m^d})X^{-1}+
((XYX^{-1})^{m^d}-(Y^m)^{m^d})$. The first term is at most $b_d+2e$ on $\xi$. For the second, the relator gives
$\hs{XYX^{-1}\xi-\xi Y^m}\le(3m+1)e$, so by~\eqref{eq:left-right-power} and $c(Y^{mn})\le mne$ it is at most $(4m+1)m^de$. Hence
$b_d\le Cm^de$, with no factor $m^{2d}$.

Conjugating by $X^{\min(j,l)}$, at cost $O((N+1)e)$, reduces $[Z_j,Z_l]$ to $[Y,X^{|l-j|}YX^{-|l-j|}]$, and the second factor is
within $Cm^{|l-j|}e$ of $Y^{m^{|l-j|}}$, which commutes with $Y$. Applying~\eqref{eq:left-right-power} to $UYU^*$ and $\zeta Y$ gives
$UY^nU^*$ within $O(ne)$ of $\zeta^nY^n$, and $\zeta^{m^d}=1$ for $d\ge1$ because $q\mid m$; comparing with $b_d$ and conjugating by
$X^j$ gives the estimate for $[T_j,Z_l]$. The conjugacy estimate is the defining property of $U$ conjugated by $X^j$.
Words of different copies have length $O(N+1)$ and commute after $O((N+1)^2)$ swaps, each with a suffix of length $O(N+1)$.
\end{proof}

For a unitary $Z$ and an offset $\vartheta$ uniform modulo $1/q$, let $P_a(Z;\vartheta)$, $a\in\Z_q$, be the spectral projections of
$Z$ onto the arcs $[\vartheta+a/q,\vartheta+(a+1)/q)$; multiplying $Z$ by $\zeta$ permutes them cyclically, for every offset. For
partitions put $C(P)=(\sum_a\hs{[P_a,\xi]}^2)^{1/2}$, $D(P,V)=(\sum_a\hs{[P_a,V]\xi}^2)^{1/2}$ and
$D(P,Q)=(\sum_{a,b}\hs{[P_a,Q_b]\xi}^2)^{1/2}$.

\begin{lemma}[Random spectral cuts]\label{lem:random-cuts}
For unitaries $R,S$, $\mathbb E\sum_a\hs{P_a(R)\xi-\xi P_a(S)}^2\le\frac q2\hs{R\xi-\xi S}$. Consequently, for unitaries $Z,V,B$,
\begin{gather*}
\mathbb EC(P(Z))^2\le C_qc(Z),\qquad\mathbb ED(P(Z),V)^2\le C_q\bigl(\hs{[Z,V]\xi}+c(Z)+c(V)\bigr),\\
\mathbb ED(P(Z),P(B))\le C\bigl(\hs{[Z,B]\xi}+c(Z)+c(B)\bigr)^{1/4},
\end{gather*}
and if $V$ conjugates $Z$ approximately to $\zeta Z$, with $\pi$ the induced cyclic permutation of the arcs,
$\mathbb E\sum_a\hs{(P_a(Z)V-VP_{\pi(a)}(Z))\xi}^2\le C_q(\hs{(VZV^*-\zeta Z)\xi}+c(Z)+c(V))$.
\end{lemma}

\begin{proof}
The weights $\hs{E_R(dt)\xi E_S(du)}^2$ form a coupling of the $\rho$-spectral laws of $R$ and $S$ with
$\int|e^{2\pi it}-e^{2\pi iu}|^2=\hs{R\xi-\xi S}^2$ (\cref{lem:spectral-coupling}). The left side of the first claim is twice the
coupling mass of pairs in different arcs, and a random offset separates $t$ from $u$ with probability at most $q\,d(t,u)$,
where $d(t,u)\le\frac14|e^{2\pi it}-e^{2\pi iu}|$; Cauchy--Schwarz gives the bound. Take $R=S=Z$ for the first consequence, and
$R=VZV^*$, $S=Z$ or $S=\zeta Z$ for the second and the covariant one, converting $P_a(R)\xi-\xi P_a(S)$ into the stated
commutators at the cost of $C(P(Z))$ and $c(V)$; all quantities are bounded, so squares may be absorbed into first powers. For
the third, the second applied with $Z$ as control bounds $\mathbb E\sum_b\hs{[Q_b,Z]\xi}^2$ for the arcs $Q_b$ of $B$; applying it again
to the arcs of $Z$ with the unitary controls $2Q_b-I$, Cauchy--Schwarz over $b$ and Jensen give the fourth root.
\end{proof}

\begin{lemma}[Weighted histories]\label{lem:histories}
Let $P_1,\dots,P_M$ be projective partitions with at most $Q$ outcomes, $c_t=C(P_t)$, $c=\sum_tc_t$, $d_{t,s}=D(P_t,P_s)$, and
let $K_h$ be the Kraus products of the sequential measurement, in physical order. Then
$\sum_hK_h^*K_h=\sum_hK_hK_h^*=I$ and, for every operator $F$,
\[
\Bigl(\sum_h\hs{[K_h,\xi]}^2\Bigr)^{1/2}\le c,\qquad\Bigl(\sum_h\hs{FK_h\xi}^2\Bigr)^{1/2}\le\hs{F\xi}+\opn Fc .
\]
Let $B$ be Bob's record of the sequence on $\rho$.
\begin{enumerate}[label=(\alph*),itemsep=1pt]
\item If a unitary $V_t$ fixes the earlier labels and permutes the $q_t$ labels of step $t$ cyclically, and $a_t$ is the sum of
$c(V_t)$, of the aggregate covariance errors $(\sum_a\hs{(P_{s,a}V_t-V_tP_{s,\pi_s(a)})\xi}^2)^{1/2}$ over $s\le t$, and of $C_QMc$,
then $H(B_t\mid B_1,\dots,B_{t-1})\ge\log q_t-C_Q\min(1,\sqrt{a_t})$.
\item On the canonical purification, let Alice measure the transposed partitions in the same order, with record $A$. Then
$H(B\mid A)\le C_Q\sum_t\min(1,c_t+\sum_{s<t}d_{t,s}+Mc)$, and $S(\rho)\ge H(B)-H(B\mid A)$.
\end{enumerate}
\end{lemma}

\begin{proof}
Completeness follows by summing the innermost factor. Expanding $[AB,\xi]=A[B,\xi]+[A,\xi]B$ and summing each term by one of
the two completeness relations gives the centrality bound by induction; then $FK_h\xi=F\xi K_h+F[K_h,\xi]$, and
$\sum_h\hs{F\xi K_h}^2=\hs{F\xi}^2$ by unitality.

(a) Telescope $V_tK_h-K_{\pi h}V_t$ through the prefix; each term is a later product, a covariance defect
$P_{s,a}V_t-V_tP_{s,\pi_s(a)}$, and an earlier history applied to $\xi$. Summing the later outcomes by completeness and the earlier
histories by the second estimate bounds the aggregate by $a_t$; moving $V_t$ through $\xi$ costs $c(V_t)$. Since right
multiplication by $V_t$ preserves the norm, the prefix law $\hs{K_h\xi}^2$ and its cyclic shift in the current label are within
total variation $Ca_t$. Averaging over the $q_t$ powers of the shift gives a law uniform in the current label given the past,
and continuity of conditional entropy over an alphabet of fixed size gives the claim.

(b) The joint components before step $t$ are $\sigma_{h,g}=K_h\xi K_g^*$, with $h$ Bob's and $g$ Alice's history. For every
matrix $\sigma$, $\sum_a\hs{P_{t,a}\sigma-\sigma P_{t,a}}^2=2\sum_{a\ne b}\hs{P_{t,a}\sigma P_{t,b}}^2$, so $\sqrt{2\Pr[A_t\ne B_t]}$ is the
aggregate norm of $[P_{t,a},\sigma_{h,g}]=[P_{t,a},K_h]\xi K_g^*+K_h[P_{t,a},\xi]K_g^*+K_h\xi[P_{t,a},K_g^*]$. The middle term
aggregates to exactly $c_t$, and the two outer terms to the aggregate of $[P_{t,a},K_h]\xi$, which telescopes to at most
$\sum_{s<t}d_{t,s}+C_QMc$. Fano's inequality with $h(x)\le C\sqrt x$ and subadditivity bound $H(B\mid A)$. Finally Alice's
measurement prepares on Bob an ensemble averaging to $\rho$, so Holevo's bound gives $I(A:B)\le S(\rho)$.
\end{proof}

\begin{proof}[Proof of \cref{thm:additive}]
Put $n_i=\lfloor\log_{m_i}(1/e)\rfloor$, $N=\max_in_i$ and $M=\sum_in_i$, so $M,N=O(\log(1/e))$ and $em_i^{n_i}\le1$. For each copy
measure the $q_i$ arcs of $Z_{i,j}$ for $j=n_i-1,\dots,0$, in this reverse order, the copies in a fixed order, with
independent random offsets; this is an actual sequence of $M$ projective measurements. By \cref{lem:short-words} and
\cref{lem:random-cuts}, $\mathbb Ec\le CM\sqrt{(N+1)e}$, so every polynomial in $M,N$ times $\mathbb Ec$ tends to zero.

\emph{Bob's record.} At stage $(i,j)$ use $V=T_{i,j}$. It permutes the current arcs cyclically, and it nearly fixes the arcs of
every earlier stage: earlier stages of the same copy are deeper, $l>j$, and those of other copies are controlled by the
cross-copy estimates. By \cref{lem:random-cuts} and a geometric sum,
$\mathbb Ea_{i,j}\le C(em_i^{n_i-j})^{1/2}+C(N+1)^Ce^{1/2}$, so by \cref{lem:histories}(a) and Jensen
$\mathbb E[\log q_i-H(B_{i,j}\mid\text{earlier})]\le C(em_i^{n_i-j})^{1/4}+C(N+1)^Ce^{1/4}$. The first terms form a geometric series
in $j$, bounded uniformly in $e$, and the remainders sum to $o(1)$. Hence $\mathbb EH(B)\ge\sum_in_i\log q_i-C$; the conditional
entropies add across copies because the earlier record contains the copies already processed.

\emph{Alice's prediction.} For two stages of one copy, $\mathbb Ed\le C(em_i^{|l-j|})^{1/4}+C((N+1)e)^{1/4}$, and for different copies
$\mathbb Ed\le C((N+1)^3e)^{1/4}$. At stage $(i,j)$ the earlier stages of the same copy have $l>j$, and
$\sum_{l>j}(em_i^{l-j})^{1/4}\le C(em_i^{n_i-j})^{1/4}$. Summing \cref{lem:histories}(b) gives $\mathbb EH(B\mid A)\le C$.

\emph{Conclusion.} For every offset $S(\rho)\ge H(B)-H(B\mid A)$; averaging over offsets, which does not change $\rho$,
$S(\rho)\ge\sum_in_i\log q_i-C=s\log(1/e)-C$. Finally $\log\rank\rho\ge S(\rho)$. For the upper bound, tensor the Fourier-lamp
models of the copies at precision $e$, with the flat state: cross-copy commutators vanish exactly, the marked involution is
$-I$ in every factor, internal defects are unchanged, dimensions multiply and entropies add.
\end{proof}

The estimates of \cref{lem:random-cuts,lem:histories} are written for this purpose, but each is elementary: the first is
Connes' coupling with a random cut, the second uses that a sequence of projective pinchings is both trace preserving and
unital, so that weighted errors can be summed over all histories without a factor of their number. The reverse order of
depths is what makes the same-copy errors geometric.

\subsection{Shared dilations force uniformity}\label{sec:rates-shared}

Two dilations acting on the same $Y$ can force more than either alone. Take commuting involutions $z_0,z_1,w$ with the row
$z_0z_1w=1$, a unitary $y$ exchanging $z_0$ and $z_1$ by conjugation, and unitaries $x_4,x_6$ with $x_jwx_j^{-1}=w$ and
$x_jyx_j^{-1}=y^j$; the relation $[x_4,x_6]=1$ may be added. Call this source $K_{4,6}$; the marked involution is $w$.

\begin{theorem}[Shared dilations $4$ and $6$]\label{thm:shared}
Every model of $K_{4,6}$ at precision $e$, with or without $[x_4,x_6]=1$, satisfies $d\ge\rank\rho\ge c/e$ and
$S(\rho)\ge\log(1/e)-C$, and there are flat models of dimension $O(1/e)$ satisfying $[x_4,x_6]=1$ exactly. For the
$\rho$-spectral law $\mu$ of $Y$, with $\delta_j=\Wone(j_*\mu,\mu)$ and $\delta_S=\Wone((+\frac12)_*\mu,\mu)$ in turns,
\[
\Wone(\mu,h)\ \le\ \tfrac13(4\delta_6+7\delta_4+5\delta_S)\ \le\ \tfrac{85}4e .
\]
\end{theorem}

\begin{proof}
\emph{Spectral defects.} The lamp swap and the row give $\hs{(Z_1+Z_0)\xi}\le5e$, hence $\hs{(Z_0YZ_0^*+Y)\xi}\le12e$ after moving
suffixes, and Connes' coupling with the state-change cost gives $\delta_S\le\frac{17}4e$. The dilation relators give
$\delta_j\le\frac{j+3}2e$, so $\delta_4\le\frac72e$ and $\delta_6\le\frac92e$.

\emph{Inverse-root kernels.} For $m\ge1$ let $D_m$ be the pushforward under $t\mapsto mt$, $Q_m$ the average of the rotations by
$j/m$, and $R_m$ the operator choosing a uniform $m$-th root. Then $R_mD_m=Q_m$, $R_2Q_2D_2=Q_4$,
$\Wone(R_m\alpha,R_m\beta)\le\frac1m\Wone(\alpha,\beta)$ (couple $t$ with $t+\Delta$ by a shortest signed difference and the roots
$(t+j)/m$ with $(t+\Delta+j)/m$), and, as $R_m$ fixes Haar measure, $\Wone(\mu,h)\le\frac m{m-1}\Wone(\mu,R_m\mu)$.

\emph{The arithmetic.} Since $D_2D_6=D_3D_4=D_{12}$, $\Wone(D_2\mu,D_3\mu)\le2\delta_6+3\delta_4$. The map $D_3$ commutes with the
half-turn because $3$ is odd and is $3$-Lipschitz, so $\Wone(D_2\mu,S_*D_2\mu)\le4\delta_6+6\delta_4+3\delta_S$. Then
$\Wone(\mu,Q_4\mu)\le\Wone(\mu,Q_2\mu)+\frac14\Wone(D_2\mu,S_*D_2\mu)\le\delta_6+\frac32\delta_4+\frac54\delta_S$, and
$\Wone(\mu,R_4\mu)\le\Wone(\mu,Q_4\mu)+\frac14\delta_4$. With $m=4$ this proves the display.

\emph{Lower bounds.} A law with at most $d$ atoms is at distance at least $1/(4d)$ from Haar measure in turns, so $d\ge1/(85e)$.
\Cref{thm:entropy-transport} with $\nu=h$ and $\hs{[Y,\xi]}\le e$ gives $S(\rho)\ge\log(1/e)-C$, and $\rank\rho\ge2^{S(\rho)}$.

\emph{Upper models.} For odd $L$ put $d=2^L$ and $N=d-1$, which is coprime to $6$; let $Y e_b=e^{2\pi ib/N}e_b$ for
$0\le b\le N$ (one duplicated root at $b=N$), $Z_0e_b=e_{b\oplus2^{L-1}}$, $Z_1=-Z_0$, $W=-I$, and $X_j$ the permutation matrix
with $X_je_b=e_{j^{-1}b\bmod N}$ on the cyclic labels $b<N$, where $j^{-1}$ is the inverse of $j$ modulo $N$, fixing the additional
label $N$; equivalently $(X_j)_{b,jb}=1$. Then $X_j^*e_b=e_{jb}$, so $X_jYX_j^*e_b=e^{2\pi ijb/N}e_b$ and $X_jYX_j^*=Y^j$; with the opposite
convention $X_je_b=e_{jb}$ one would get $Y^{j^{-1}}$ instead, for instance $Y^2$ at $N=7$, $j=4$. Also $X_jWX_j^*=W$ and $X_4X_6=X_6X_4$ exactly, and flipping the highest
bit changes $b/N$ by $\pm(\frac12+\frac1{2N})$, so the lamp swaps have operator defect at most $\pi/N$. With the flat state this is a
model at precision $\pi/N$ and entropy $\log d$.
\end{proof}

Each dilation alone permits the smaller exponent $\log_42=\frac12$ or $\log_62$; sharing $Y$ forces uniformity. This is an exact
elementary transport identity, not a quantitative solution of the $\times m,\times n$ rigidity problem: for the pair $4,9$ the
present bounds are $\Omega(e^{-1/2})$ and $O(e^{-1})$ (\cref{sec:open}).

